\section{Experimental design}
\label{sec:experimental-design}

The experiments ask two kinds of question. The first asks whether the reconciliations that exist hold when they are run again: the printed tables, the independent routes, the closed forms and the Monte Carlo reference. The second asks what the acceptance rule of the Monte Carlo reconciliation is worth, which no committed document measures: its size, its power against defects of known size, and the contribution of each of its parts. Every number comes from a script in the \texttt{experiments} folder of the manuscript, which reads the committed documents of the product or calls the engine \texttt{riskvalidation} 1.1.0 and the product's generators, with the seeds and repetitions stated below.

\subsection{Hypotheses and predictions}
\label{sec:design-hypotheses}

The hypotheses are the claims that the product's design document makes for its reconciliations, quoted from it; the document sets no numerical refutation threshold beyond the claims themselves, so each claim is refuted by its negation. They restate evidence that the product had committed before the experiments of this work were designed, so they are tests of whether committed claims reproduce, and not pre-registered predictions. Table~\ref{tab:hypotheses} lists them.

\begin{table}[htbp]
\centering
\small
\caption{The six hypotheses, quoted from the product's design document, with the reference class of each and the experiment that tests it; the design document fixes no numerical refutation threshold, so each hypothesis is refuted by the negation of its claim.}
\label{tab:hypotheses}
\begin{tabularx}{\textwidth}{@{}l>{\raggedright\arraybackslash}X l l@{}}
\toprule
ID & Claim as quoted & Class & Tested by \\
\midrule
H1 & ``the printed tables and papers recomputed beside their prints (413 printed numbers of five papers in the ratings and calibration cases, 296 reproduced to the printed digit, every other one listed with its cell and its explanation)'' & P & \texttt{exp\_printed} \\
H2 & ``the p-values of the Jeffreys and binomial tests against two recomputations that call neither of the engine's SciPy functions (a quadrature of the Beta density and an exact rational sum of the binomial terms), equal within 1e-9 on a grid of 36 cases'' & I & \texttt{exp\_routes} \\
H3 & ``the Hull-White exposure of the single swap and of the four-swap netting set against ORE 1.8.17.0, within two combined Monte Carlo standard errors at all 33 report dates, with the EEPE and the CVA'' & M & \texttt{exp\_committed}, \texttt{exp\_ore\_replicates} \\
H4 & ``the Hull-White closed forms against QuantLib 1.44 to a relative 1e-6'' & F & \texttt{exp\_closedforms} \\
H5 & ``the wrong-way calibration against the printed values of the Appendix of Hull and White (2012)'' & P & \texttt{exp\_printed} \\
H6 & ``the collateralised variants (CSA with 10 and with 20 days) and the initial-margin variant against ORE [reported as gaps], with the largest EE gaps of 9.9, 7.1 and 11.7 combined standard errors, CVA gaps of 15.9 and 9.9 in the two CSA variants, and a mean dynamic initial margin of ORE at 0.88 of the engine's'' & M & \texttt{exp\_committed}, \texttt{exp\_ore\_replicates} \\
\bottomrule
\end{tabularx}
\end{table}

Five predictions about the repeated engine-against-ORE comparison were written by the author in a file of the experiments folder before the first replicate was run, and are reported as held or not held in Section~\ref{sec:results}. R2 was written after a pilot of the size-and-power experiment had measured a spread near $0.6$ for two runs of the engine; it therefore tests whether that spread transfers to a pair made of the engine and ORE, and it is not a blind prediction. The repository history cannot show the order of the predictions and the runs, because the file and the outputs enter in one commit; only the modification times of the files, which put the predictions first, record it. They state that the mean standardised difference at each date of the uncollateralised variants lies within $0.5$ of zero and the mean over the dates of those means within $0.25$ (R1); that the spread of that difference over repetitions is below $0.8$ with the product's standard errors and between $0.8$ and $1.2$ with pair-aware ones (R2); that the product's rule passes at least half of the repetitions (R3); that the rule fails in every repetition of the 10-day collateralised variant with a mean largest $\lvert z\rvert$ above $5$ (R4); and that the valuation-adjustment difference of that variant is below $-5$ in at least $7$ of $8$ repetitions (R5). The experiment on the size and power of the rule was designed after a pilot of four repetitions that showed the spread of the standardised difference to be well below one; no prediction about it was fixed in advance, and its results are reported as measurements. Three predictions about the maximum-T comparator of Section~\ref{sec:design-methods} (M1 to M3) were written in the same file after the review of the first version of this work and before that comparator was run for the first time: that the size of the two maximum-T rules lies between $3\,\%$ and $8\,\%$ (M1); that the one with pair-aware errors detects a $2\,\%$ volatility error at least as often as the pair-aware rule at the bound $3.18$ (M2); and that the one with the product's errors detects it less often than the one with pair-aware errors (M3).

\subsection{Data and provenance}
\label{sec:design-data}

\textbf{Printed answers.} The five papers are Pluto and Tasche \citep{plutotasche2005} (Tables 1 to 14: the most prudent default probability of three grades of 100, 400 and 300 obligors, with 0 defaults or with 0, 2 and 1), Tasche \citep{tasche2013} (Tables 5 to 9, from the rates and profiles of the paper's example for 2009, 2010 and 2011), Schuermann and Hanson \citep{schuermann2004} (Table 5, the Wald and Agresti--Coull intervals of 531 obligors), Israel, Rosenthal and Wei \citep{israel2001} (nine $L_1$ distances between generators) and Engelmann \citep{engelmann2024} (the through-the-cycle portfolio and the projected default probabilities). All are sources of the derived-only licence class: the committed documents hold the printed numbers, rates and profiles that the engine needs and never the papers' own tables of counts, which is why the inputs of the sources read as recorded (Israel et al., Engelmann, and Tables 13 and 14 of Pluto and Tasche) are not recomputed here and their gaps are read as the documents record them.

\textbf{The exposure case.} The swap book is synthetic: a 10-year payer swap of 10 million USD, a 5-year receiver of 6 million, a 3-year payer of 4 million and a 6-month receiver of 2 million, with fixed rates equal to the par yield at each maturity on 30 September 2025 rounded to 10 basis points. The curve is the United States Treasury daily par yield curve of that date (licence class mirror-allowed), converted to continuously compounded zero rates, linear in the rate. The Hull--White parameters are estimated by the exact autoregressive least-squares fit of the stochastic-processes toolkit of Brigo and co-authors \citep{brigo2008} on the 0.25-year Treasury zero rate from 2 January 1990 to 29 September 2025 (\nHWobs\ observations): mean reversion $\nHWa$ with standard error $\nHWaSe$ and volatility $\nHWsigma$ with standard error $\nHWsigmaSe$. The 33 report dates are monthly to one year, quarterly to three years and half-yearly to nine and a half years. The counterparty curve is the spread of the Treasury high quality market corporate bond curve over the Treasury par curve of the same month end, with the credit triangle $S(t)=\exp(-st/\mathrm{LGD})$ and a loss given default of $0.6$. The collateral terms are a zero threshold, a minimum transfer amount of 10,000 USD, margin periods of risk of 10 and 20 business days and a dynamic initial margin by regression at the $99\,\%$ quantile. The exposure case is the product's own and its inputs are rebuilt here from the committed documents.

\textbf{Software.} The engine is \texttt{riskvalidation} 1.1.0 with SciPy 1.18.1 and NumPy 2.4.6, the Monte Carlo reference is the Open Source Risk Engine 1.8.17.0 run natively from its Windows wheel, and the closed-form library is QuantLib 1.44. The Monte Carlo reference uses a linear Gauss--Markov model with the same mean reversion and volatility, $10{,}000$ samples with antithetic pairs; the engine uses $20{,}000$ paths with antithetic pairs.

\subsection{Reference methods and comparators}
\label{sec:design-methods}

The published methods compared are the references themselves, all run by the product or by this work's scripts on the same inputs as the engine and none taken from a reported number: the five papers' formulas as implemented by the engine and recomputed (Algorithm~\ref{alg:printed}); the quadrature and the exact sum for the two $p$-values; QuantLib's Hull--White bond option and Jamshidian swaption; ORE's exposure simulation; and the engine's own analytic identity for the expected exposure of a single swap at its reset dates, which the product runs as a test with a gate of three standard errors. The engine's cohort estimator has a comparator the engine carries, the \texttt{transitionMatrix} library; its comparison is run in Section~\ref{sec:results} on a seeded panel.

For the acceptance rule of the Monte Carlo reconciliation the comparators are twelve rules applied to the same pairs of runs: the product's rule (all $33$ dates, the EEPE and the CVA within two combined errors); its three components alone (the dates, the EEPE, the CVA); the product's rule at the family-wise bound $c_{35}$; the same rule with pair-aware errors, at the bound $2$ and at $c_{35}$; a tolerance on the largest relative difference of the profile, with no standard error, at $1\,\%$, $2\,\%$ and $5\,\%$ of the peak; and the maximum-T bound, which keeps the family-wise size under dependence: the largest of the 35 absolute standardised differences, with the product's errors or with pair-aware errors, against the $95\,\%$ point of the largest absolute value of a Gaussian vector with the covariance of the 35 differences under the null \citep{hothorn2008}. That covariance is estimated on the null repetitions of the other half of the repetitions, so that no repetition is judged with a bound it helped to set. It comes from repetitions of the engine against itself, which a validation function can run as pilots; it is not estimated from the single pair of runs of a reconciliation, and the maximum-T rows of Section~\ref{sec:results-power} are therefore the size and power that a dependence-aware bound can reach when its covariance is known, not what a reconciliation against ORE can compute from its own two runs. The global $\chi^{2}$ test of the whole profile is not run: it answers whether the profile differs and not at which date or functional, which the reconciliation reports, and its reference distribution needs the same covariance, of rank at most the number of repetitions of the half it is estimated on.

\subsection{Metrics}
\label{sec:design-metrics}

The metrics are those of Section~\ref{sec:proposed-method}: the digit rule~\eqref{eq:digit} and the gap in units~\eqref{eq:unitgap} for printed answers; the maximum absolute difference~\eqref{eq:route-tol} for independent routes; the relative difference~\eqref{eq:closed-rel} for closed forms; the standardised difference~\eqref{eq:z} and the rule~\eqref{eq:rule} for ORE. The size and the power of a rule are rejection rates over repetitions, each with its Wilson $95\,\%$ interval; the pair correlation $\rho$ is the Pearson correlation between the discounted exposure of path $i$ and of path $i+P/2$ at a date; and the equivalence margin and the smallest detectable gap are Equation~\eqref{eq:equivalence} and Proposition~\ref{prop:power}. The Wilson interval of a rejection rate $\hat{p}=k/R$ over $R$ repetitions, with $z=1.96$, is
\begin{equation}
\frac{\hat{p}+z^{2}/(2R)\ \pm\ z\sqrt{\hat{p}(1-\hat{p})/R+z^{2}/(4R^{2})}}{1+z^{2}/R},
\label{eq:wilson}
\end{equation}
and the correlation between the two paths of a pair at date $t_k$ is estimated over the $P/2$ pairs by
\begin{equation}
\hat\rho_k\ =\ \frac{\sum_i\bigl(x_{ik}-\bar{x}_{k}\bigr)\bigl(x_{i+P/2,k}-\bar{x}'_{k}\bigr)}{\sqrt{\sum_i\bigl(x_{ik}-\bar{x}_{k}\bigr)^{2}\ \sum_i\bigl(x_{i+P/2,k}-\bar{x}'_{k}\bigr)^{2}}},
\label{eq:rho-hat}
\end{equation}
with $x_{ik}$ the discounted exposure of path $i$ at date $k$ and the bars the means over the first and the second half of the paths. The spread of the standardised difference at a date is the standard deviation of $z_k$ over the repetitions, and the proposition on antithetic pairs predicts it as $\sqrt{1+\bar\rho_k}$ from $\hat\rho_k$ averaged over simulations. The scalar functionals are compared in the same units: the EEPE with its batch error, the CVA with its path error, and the equivalence margin and detectable gap as shares of the engine's value.

\subsection{Protocol, repetitions and seeds}
\label{sec:design-protocol}

The printed-answer reconciliation is deterministic. The grid of the independent routes is the product's grid of $36$ cases, extended by $\nRandomCases$ cases drawn with seed $20261104$ ($n$ log-uniform on $[10,5000]$, $D$ uniform on $0$ to $\min(n,50)$, $x$ log-uniform on $[10^{-4},0.3]$). The closed forms are compared on the four cases of the product and on $\nClosedCases$ cases of a seed-$20261104$ draw of mean reversion, volatility, expiry, maturity and strike or fixed rate. The planted defects are a tail convention moved by one in the binomial helper and a relative error $\epsilon$ from $10^{-12}$ to $10^{-3}$ inside the SciPy functions.

The size and power of the rules are measured with $\nCritReps$ repetitions. In repetition $r$ the reference is a run of the engine of $\nCritPathsB$ paths with seed $1{,}000{,}000+2r+1$ at the calibrated parameters; the engine under test is a run of $\nCritPathsA$ paths with seed $1{,}000{,}000+2r$ under seven conditions: the calibrated parameters (the null), the volatility raised by $2$, $5$ and $10\,\%$, the mean reversion raised by $50$ and $100\,\%$, and a measure error that reports the discounted expected exposure where the forward-measure one is meant. The truth is known, so every rejection of the null is a false alarm and every acceptance of a defect a missed detection. The repeated comparison with ORE uses $\nOreRepsNet$ repetitions for each of the two uncollateralised variants (seeds $5000$ to $5023$, the same seed for the engine and for ORE, whose generators differ) and $\nOreRepsCsa$ for the 10-day collateralised variant with the diagnostic run of zero minimum transfer amount, and a single repetition at the product's seed $42$ for each variant, to check the committed run. The pair correlation is measured on $\nPairReps$ simulations with seeds $2{,}000{,}000+r$.

\subsection{Statistics}
\label{sec:design-statistics}

A rejection rate is reported with its Wilson interval \citep{brown2001}. A mean of standardised differences over repetitions is reported with its standard deviation over repetitions, from which its standard error follows; the standard deviation over repetitions is the quantity that the calibration of the standard error is judged by (Proposition~\ref{prop:antithetic}). The comparison of rules is made on the same pairs of runs, so the difference between two rules is a paired comparison, and the intervals of two rules that overlap are not read as equal. No hypothesis test is applied to the outcomes of the hypotheses of Table~\ref{tab:hypotheses}: each is a claim about a count or a bound that the measurement meets or does not.

The repetitions of one experiment are independent: each has its own seeds, and no random number is shared between two repetitions. Within a repetition the conditions of the size and power experiment share the seed of the engine's run and the run of the reference, which are common random numbers; the difference between two conditions in one repetition is therefore a paired difference, and the rejection rates of two conditions are not independent estimates. With $R=\nCritReps$ repetitions the Wilson half-width of a rate near $0.5$ is about $0.05$, and that of a rate near $0.05$ about $0.02$, which is the resolution of every size and power in the tables; a difference of two sizes smaller than that is not read. The standard deviation over repetitions is reported with the number of repetitions it rests on, and no interval is put on it, because it is used only to compare with a value of one that is far from the measured spread. The independent routes and the closed forms are deterministic, so their results are single numbers, with the seeds of their parameter draws fixed and stated.
